Table~\ref{matrix_xg_train} shows confusion matrices of five XGBoost models: four with a single feature and one with all features.
Table~\ref{result_xg_train} shows the evaluation metrics.
The all-feature model achieved the highest accuracy.
The model with centroid 90th percentile values achieved the lowest FNR.
Table~\ref{feature_importance_xg} shows feature importance from the all-feature model (Table~\ref{matrix_xg_train_e}).
Centroid-50th and Centroid-90th showed the highest importance.
Centroid features reduced misclassifications among similar activity levels.
A diverse set of three-dimensional features is effective for activity level estimation with XGBoost on short video datasets.
For the XGBoost model, the all-feature input outperformed every single-feature model in Table~\ref{result_xg_train}.
The contrast between RF and XGBoost indicates that complementary 3D features require an appropriate ensemble model to be fully exploited.
